\chapter{Конструкторский раздел}
\label{cha:design}

В данном разделе проектируется новая всячина.

\section{Архитектура всячины}

\paragraph{Проверка} параграфа. Вроде работает.
\paragraph{Вторая проверка} параграфа. Опять работает.

Вот.

\begin{itemize}
\item Это список с <<палочками>>.
\item Хотя он и не по ГОСТ, кажется.
\end{itemize}

\begin{enumerate}
\item Поэтому для списка, начинающегося с заглавной буквы, лучше список с цифрами.
\end{enumerate}

Формула \ref{F:F1} совершено бессмысленна.

%Кстати, при каких-то условиях <<удавалось>> получить двойный скобки вокруг номеров формул. Вопрос исследуется.

\begin{equation}
a= cb
\label{F:F1}
\end{equation}
f72{z8VdeeHz}x6Hq6Z6P82V23

Окружение \texttt{cases} опять работает (см. \ref{F:F2}), спасибо И. Короткову за исправления..


\begin{equation}
a= \begin{cases}
 3x + 5y + z, \mbox{если хорошо} \\
 7x - 2y + 4z, \mbox{если плохо}\\
 -6x + 3y + 2z, \mbox{если совсем плохо}
\end{cases}
\label{F:F2}
\end{equation}

\section{Подсистема всякой ерунды}

Культурная вставка dot-файлов через утилиту dot2tex (рис.~\ref{fig:fig02}).

\begin{figure}
  \centering
% [width=0.5\textwidth] --- регулировка ширины картинки
  % \includegraphics{inc/dot/cow2}
  \caption{Рисунок}
  \label{fig:fig02}
\end{figure}


% \begin{figure}[ht!]
%  \centering 
%  \includegraphics[width=\textwidth]{inc/raster/sample.jpg} 
%  \caption{A simple caption} 
%  \label{overflow} 
% \end{figure}



\section{Какой-то последний раздел}
Проблемы, которые придется нам решить в ходе данной дипломной работы это малое количество flash-памяти на роботе (256Кб) и недостаточная разрешающая способность блока с камерой (CAM-NXT).



\subsection{Блок-схема всякой ерунды}

\subsubsection*{Кстати о заголовках}

У нас есть и \Code{subsubsection}. Только лучше её не нумеровать.


%%% Local Variables:
%%% mode: latex
%%% TeX-master: "rpz"
%%% End:
